\selectlanguage{ngerman}%
\frenchspacing%
\begin{ledgroupsized}[r]{120mm}%
\footnotesize%
\pstart%
\noindent%
\textbf{Überlieferung:}%
\pend%
\end{ledgroupsized}%
\begin{ledgroupsized}[r]{114mm}%
\footnotesize%
\pstart%
\parindent -6mm%
\makebox[6mm][l]{\textit{L}}%
Aufzeichnung:
LH XXXV 14,\;2 Bl.\;165.
Ein Zettel (9,5 x 10,4 cm);
Ränder beschnitten;
römisches Papier (ohne Wasserzeichen).
Zwei Seiten auf Bl.\;\href{\Scan{LH_35_14_02_165v}}{165~v\textsuperscript{o}}\! und \href{\Scan{LH_35_14_02_165r}}{165~r\textsuperscript{o}}.
Auf Bl.\;\href{\Scan{LH_35_14_02_165r}}{165~r\textsuperscript{o}}\! ist quer zur Schreibrichtung das durchgestrichene Wort \textit{semper} zu lesen.
\pend%
\end{ledgroupsized}%
\vspace{5mm}%
\begin{ledgroup}%
\footnotesize%
\pstart%
\noindent%
\textbf{Datierungsgründe:}
Die Datierung der Aufzeichnung N.\;\ref{58258} beruht auf Eigenschaften des Textträgers sowie 
inhaltlichen Betrachtungen.
Siehe die ausführliche Begründung auf S.\;\pageref{58258_dat}.%
\pend%
\end{ledgroup}%
\selectlanguage{latin}%
\frenchspacing%
\vspace{8mm}%
\normalsize%
\pstart%
\noindent%
\edtext{%
\lbrack%
\href{\Scan{LH_35_14_02_165v}}%
{165~v\textsuperscript{o}}%
\rbrack\
In omni mobili%
\indexS{mobile}
partem accipere licet
cujus volumen%
\indexS{volumen mobilis}
sit ad volumen mobilis %
in ratione data%
\indexS{ratio data}
minoris ad majus,%
}{%
\lemma{\lbrack\href{\Scan{LH_35_14_02_165v}}{165~v\textsuperscript{o}}\rbrack}%
\Bfootnote{%
\textit{(1)}~In quacunque mole partem accipere licet, et rat %
\textit{(2)}~In omni %
\lbrack...\rbrack\ accipere licet %
\textit{(a)}~dato volumini minori %
\textit{(b)}~volumine mobilis minori %
\textit{(c)}~cujus volumen %
\lbrack...\rbrack\ ad majus,%
~\textit{L}%
}}
et densitas%
\indexS{densitas mobilis}
sit densitati mobilis aequalis.
\pend%
\begin{edfigure}[1.750em][1.00em]
\includegraphics%
{gesamttex/edit_VIII-4/images/LH_35_14_02_165_d_165v.pdf}
\edcaption[0.5em]{%
\lbrack\textit{Fig.\:1}\rbrack%
}%
\end{edfigure}%
\pstart%
\noindent%
Cum in toto quovis pars capi possit %
quae sit ad totum
in ratione data minoris ad majus,%
\indexS{ratio data}
poterit
\edtext{%
et in volumine dato \textit{AB}%
\indexS{volumen mobilis}
capi pars talis.
Quod si jam totum est constantis ubique densitatis,%
\indexS{densitas constans}
etiam pars erit ejusdem densitatis,
et habetur propositum.%
\indexS{propositum}
Si totum sit minoris densitatis %
quam pars haec \textit{C},
necesse est
totum idem alibi majoris esse densitatis
quam pars \textit{C};%
}{%
\lemma{et}%
\Bfootnote{%
\textit{(1)}~capi talis pars voluminis dati %
\textit{(2)}~in volumine %
\lbrack...\rbrack\ pars talis. %
\textit{(a)}~Haec pars \textit{C} %
\textit{(aa)}~aut major est toto aut minor %
\textit{(bb)}~aut densior est aut ra %
\textit{(cc)}~aut aequalis est %
\textit{(b)}~Quod si %
\lbrack...\rbrack\ propositum. Si %
\textit{(aa)}~pars %
\textit{(bb)}~totum sit %
\textit{(aaa)}~majoris %
\textit{(bbb)}~minoris densitatis, quam %
\textit{(aaaa)}~totum %
\textit{(bbbb)}~pars haec \textit{C}, necesse est totum %
\textit{(aaaaa)}~in aliis partibus majoris esse densitatis; %
\textit{(bbbbb)}~idem alibi majoris esse densitatis %
\textit{(aaaaa-a)}~, ut fiat compensatio %
\textit{(bbbbb-b)}~quam pars \textit{C};%
~\textit{L}%
}}
ibi sumatur \textit{E} aequalis ipsi \textit{C},
et transferendo \textit{C} in \textit{E},
necesse est
\edtext{%
\lbrack\textit{C}\rbrack%
}{%
\lemma{\textit{C}}%
\Bfootnote{%
\textit{erg.\;Hrsg.}%
}}
continue accedere ad densitatem ipsius \textit{E},
ergo transire per intermedias,%
\indexS{densitas intermedia}
adeoque et eam
qua aequale sit ipsi volumini \textit{AB}.%
\indexS{volumen mobilis}
Haec propter exemplum ratiocinandi.%
\indexS{exemplum ratiocinandi}
Unde jam deduco generaliter mobilia%
\indexS{mobile}
\lbrack%
\href{\Scan{LH_35_14_02_165r}}%
{165~r\textsuperscript{o}}%
\rbrack\
habere moles%
\indexS{moles mobilis}
in ratione composita
\edtext{%
\lbrack volumin\rbrack um%
}{%
\lemma{mobilium}%
\Bfootnote{%
\textit{L\;ändert Hrsg.}%
}}
et densitatum,%
\indexS{densitas mobilis}
etiamsi sint inconstantis%
\indexS{densitas inconstans}
\edtext{%
densitatis.
Ex eo
quod
\edtext{%
assumsimus
mobilia aequalis voluminis%
\indexS{mobilia voluminis aequalis}
esse in ratione densitatum,%
\indexS{densitas mobilis}
et aequalis densitatis%
\indexS{mobilia densitatis aequalis}
in ratione voluminum.%
\indexS{volumen mobilis}%
}{%
\lemma{assumsimus \lbrack...\rbrack\ voluminum}%
\Cfootnote{%
Siehe \textsc{G.\;W.\;Leibniz},
\textit{Dynamica},
pars~I, sectio~I, cap.\;II,
prop.\;1--2
(%
\textit{CFP}~I, S.\;31\:f.\;[A]; II.1, S.\;38~[C]).\cite{01354}%
}}
Sumendo nempe%
}{%
\lemma{densitatis.}%
\Bfootnote{%
\textit{(1)}~Sumendo scilicet %
\textit{(2)}~Ex eo quod assumsimus mobilia %
\textit{(a)}~esse %
\textit{(b)}~aequalis voluminis %
\lbrack...\rbrack\ Sumendo nempe%
~\textit{L}%
}}
partem majoris voluminis aequalem volumini minori,
sed aequalem densitate sui toti.
\pend%
\pstart%
Idem tamen alia ratione demonstravi
non usus hoc auxilio praesenti,%
\indexS{auxilium}
\edtext{%
in prop.\;3
mei capituli de quantitate%
\indexS{quantitas materiae}%
}{%
\lemma{in \lbrack...\rbrack\ Materiae}%
\Cfootnote{%
Vgl.\;a.a.O.,
prop.\;3
(%
\textit{CFP}~I, S.\;32\:f.\;[A]; II.1, S.\;38--40~[C]).\cite{01354}%
}}
\edtext{%
Materiae,%
\indexS{materia}
singulari illic usus artificio.%
\indexS{artificium}
Idem et alia%
}{%
\lemma{Materiae,}%
\Bfootnote{%
\textit{(1)}~idem adhuc al %
\textit{(2)}~singulari illic %
\lbrack...\rbrack\ et alia%
~\textit{L}%
}}
ratione assequi potuissem
reducendo prius mobile
ad aliud aequidensum%
\indexS{mobile aequidensum}
ejusdem molis.%
\indexS{moles mobilis}
\pend%
\vspace{3.5em}%
\pstart%
\noindent%
\centering%
*     *     *
\pend%
\vspace{3.0em}%
\selectlanguage{german}%
\frenchspacing%
\footnotesize%
\pstart%
\noindent%
\textls{Ausführliche Datierungsgründe:}%
\pagelabel{58258_dat}
Der Träger der titellosen und undatierten Aufzeichnung N.\;\ref{58258}, die mit dem Verhältnis von Dichte und Volumen eines Körpers befasst ist, weist zwar kein Wasserzeichen auf, wohl aber gleiche sichtbare Eigenschaften wie eine Papiersorte, die Leibniz während seines Aufenthalts in Rom\indexO{Rom} verwendete, unter anderem für die erste Fassung seiner \textit{Dynamica}.\cite{01354}
Auf gleichem Papier sind ebenfalls die im vorliegenden Band veröffentlichten Aufzeichnungen N.\;\ref{58656-A}, N.\;\ref{58218} und N.\;\ref{58235} überliefert %
(siehe die ausführlichen Datierungsgründe zu N.\;\ref{58656-A}, S.\;\pageref{58656-A_dat}).
Bei all diesen Texten, von denen N.\;\ref{58656-A} und N.\;\ref{58235} auch einen inhaltlichen Zusammenhang mit Leibnizens \textit{Dynamica} aufweisen, %
ist demnach eine Entstehungszeit anzunehmen, die sich auf die Monate vom Sommer bis zum November 1689 erstreckt, als Leibniz in Rom\indexO{Rom} am \textit{Phoranomus}\cite{01402} und an der ersten Fassung der \textit{Dynamica}\cite{01354} gearbeitet hat (siehe \textsc{Robinet} 1988, S.\;81\,ff.).\cite{01240}
\pend%
\pstart%
Auch Folgendes ist zu beachten:
In der Aufzeichnung N.\;\ref{58258} verweist Leibniz auf
\textit{Dynamica}, pars~I, sectio~I,
cap.\;II: \textit{De quantitate materiae seu volumine et densitate},
prop.\;3: \textit{Moles mobilium, vel ipsa mobilia, sunt in ratione composita voluminum et densitatum, seu extensionum et intensionum materiae}
(%
\textit{CFP}~I, S.\;32\:f.\;[A]; II.1, S.\;38--40~[C]).\cite{01354} %
In der ersten Fassung der \textit{Dynamica}\cite{01354} lag dieses Kapitel somit vor, als die Aufzeichnung N.\;\ref{58258} angefertigt wurde.
Dies schließt aus, dass N.\;\ref{58258} älter als die zwischen Sommer und November 1689 entstandene erste Fassung der \textit{Dynamica}\cite{01354} ist.
Eine spätere Entstehung von N.\;\ref{58258} (höchstens bis Ende März 1690, als Leibniz Italien\indexO{Italien} verlassen hat) ist zwar nicht ausgeschlossen, aber unwahrscheinlich.
\pend%
\selectlanguage{latin}%
\frenchspacing%
\newpage%
